% Modernised reading — fonds Alexandre Grothendieck, shelfmark 66, pages 1-4.
%
% This is an INTERPRETATION, not a transcription. It re-reads the pages in
% today's notation and vocabulary, states the hypotheses the manuscript leaves
% implicit, and drops the critical apparatus so that the mathematics reads
% straight through. Everything the brackets and underlines carried moves into a
% footnote.
%
% It is held to being mathematically correct as it stands, not merely faithful
% to the page. In any dispute of substance the transcription governs: whoever
% cites Grothendieck must cite batch-01.fr.tex.
%
% Pass: Opus 5 (claude-opus-5), 2026-08-16 — first pass, unchecked against the
% pages by a human. The folder was read whole; it holds one batch, pages 1-4,
% so this file covers it entire.
%
% A word on the hardest judgement in this file. The manuscript's own margins
% record what became of several subjects, and the temptation is to complete
% those records from what is now known. That is done only where the outcome is
% not in doubt and where the margin itself points at it; everywhere else the
% subject is named in modern terms and its fate left open. Saying "this became
% X" is a claim about the literature, and this edition does not make such claims
% on Grothendieck's behalf.
\documentclass[11pt,a4paper]{article}
% Le préambule commun à toutes les transcriptions du fonds Grothendieck.
%
% Il tient en une idée : **l'incertitude se compose, elle ne se résout pas.**
% Une transcription qui lisse un mot douteux a détruit la seule chose pour
% laquelle on la faisait — un lecteur doit pouvoir savoir, à chaque endroit,
% ce qui était sur le feuillet et ce qui a été ajouté. Les six macros qui
% suivent sont donc l'appareil critique, et non de la décoration.
%
% Les mêmes commandes sont reconnues par `scripts/render.mjs`, qui produit la
% vue de lecture du site. Une macro ajoutée ici doit l'être aussi là-bas,
% faute de quoi le rendu échouera bruyamment — ce qui est voulu : un rendu
% silencieusement faux est pire qu'un rendu absent.

\ProvidesPackage{grothendieck}[2026/08/08 Transcriptions du fonds Grothendieck]

\usepackage{amsmath,amssymb,amsthm}
\usepackage{mathrsfs}
\usepackage{stmaryrd}
% Les diagrammes commutatifs sont partout dans ce fonds : c'est la langue
% même de la géométrie algébrique de Grothendieck, et une transcription qui
% les rendrait en prose ne transcrirait rien.
\usepackage{tikz-cd}
% Français par défaut — c'est la langue des feuillets — mais l'anglais est
% chargé aussi : la traduction et le résumé sont des documents anglais, et la
% typographie française (espaces avant les deux-points) y serait une faute.
% Ces documents basculent par \selectlanguage{english} en tête de corps.
\usepackage[english,french]{babel}
\usepackage{fontspec}
\usepackage{geometry}
\geometry{a4paper,margin=25mm,marginparwidth=28mm}
\usepackage{marginnote}
\usepackage{xcolor}
% ulem, not soul. `soul`'s \st and \ul are built on a character-by-character
% scan that XeLaTeX's font machinery breaks: the document fails with an
% undefined control sequence rather than degrading. ulem does the same two jobs
% and is Unicode-engine safe. `normalem` stops it hijacking \emph, which the
% transcriptions use for Grothendieck's own emphasis.
\usepackage[normalem]{ulem}
\usepackage{hyperref}
\hypersetup{colorlinks=true,linkcolor=black,urlcolor=black,pdfborder={0 0 0}}

% --- Métadonnées du lot -----------------------------------------------------
% Reprises telles quelles par le script de rendu, qui les affiche en tête de
% la vue de lecture. Elles fixent ce que le fichier prétend couvrir : sans
% elles, une transcription flotte sans point d'attache dans un fonds de seize
% mille feuillets.

\newcommand{\folder}[1]{\gdef\grothFolder{#1}}
\newcommand{\batch}[1]{\gdef\grothBatch{#1}}
\newcommand{\leaves}[2]{\gdef\grothFirst{#1}\gdef\grothLast{#2}}
\newcommand{\foldertitle}[1]{\gdef\grothFoldertitle{#1}}
\gdef\grothFolder{?}\gdef\grothBatch{?}\gdef\grothFirst{?}\gdef\grothLast{?}\gdef\grothFoldertitle{}

% --- Le résumé --------------------------------------------------------------
%
% Il ouvre la lecture modernisée, sous le titre « Résumé », et il est composé
% en petit. Ce n'est pas un ornement : c'est le seul passage du document qui
% ne soit pas des mathématiques, et un lecteur doit pouvoir le reconnaître
% avant de l'avoir lu — puis le sauter s'il connaît déjà le sujet, ou s'y
% arrêter s'il ne le connaît pas. Le filet qui le suit marque la reprise du
% corps.

\newenvironment{resume}
  {\small\setlength{\parskip}{0.45em}}
  {\par\medskip\hrule\medskip}

% --- Mots-clés ---------------------------------------------------------------
%
% En anglais, en clôture du résumé de la lecture modernisée : le vocabulaire
% moderne sous lequel chercher ce que les feuillets construisent. Le manifeste
% du site (`npm run manifest`) les extrait pour étiqueter la cote entière —
% cette ligne est la source unique des tags, il n'y a pas de fichier à côté.
\newcommand{\keywords}[1]{\par\smallskip\noindent
  {\footnotesize\textsc{Keywords} --- \textit{#1}}\par}

% --- L'appareil critique ----------------------------------------------------

% Le numéro de feuillet, en marge. C'est l'ancre qui permet de régler un
% désaccord sur une formule en regardant *un* feuillet plutôt qu'une chemise
% de six cents. Le site s'en sert aussi pour tourner les pages du fac-similé
% quand on fait défiler la transcription.
\newcommand{\leaf}[1]{%
  \marginnote{\footnotesize\color{gray!70}\textsf{#1}}%
  \label{leaf:#1}\ignorespaces}

% Illisible : on ne devine pas. Un mot inventé qui se lit comme les autres est
% le pire résultat possible.
\newcommand{\ill}{\textcolor{red!60!black}{[\,\ldots\,]}}

% Lecture douteuse : le mot est donné, et signalé comme incertain.
\newcommand{\uncertain}[1]{\uline{#1}}

% Ajout de l'éditeur — une lettre manquante, un mot sous-entendu.
\newcommand{\add}[1]{\textcolor{blue!50!black}{[#1]}}

% Biffé par Grothendieck lui-même. Ce qu'il a rayé fait partie du document :
% une rature dit souvent où la pensée a changé de direction.
\newcommand{\struck}[1]{\sout{#1}}

% Note du transcripteur, sur l'état matériel du feuillet.
\newcommand{\note}[1]{\par\smallskip{\small\color{gray!60!black}\textit{[#1]}}\par\smallskip}

% Note marginale de Grothendieck, distincte du corps du texte.
\newcommand{\marginal}[1]{\par\smallskip\noindent\hspace{1em}%
  {\small\color{gray!60!black}#1}\par\smallskip}

% --- Titre ------------------------------------------------------------------

\AtBeginDocument{%
  \begin{center}
    {\large\bfseries Fonds Alexandre Grothendieck --- cote n\textsuperscript{o}~\grothFolder}\\[2pt]
    {\itshape\grothFoldertitle}\\[4pt]
    {\small feuillets \grothFirst--\grothLast\ (lot \grothBatch)}\\[2pt]
    {\footnotesize\color{gray!60!black}Transcription établie d'après le fac-similé numérisé
     par l'Université de Montpellier.\\
     Les lectures incertaines sont soulignées, les illisibles notées [\,\ldots\,],
     les ajouts entre crochets.}
  \end{center}
  \vspace{1em}}

\endinput


\folder{66}
\pages{1}{4}
\dating{1964}
\watermark{Édition de démonstration\\interprétation personnelle de l'œuvre}
\foldertitle{Problèmes : notes manuscrites (1964) — lecture modernisée du dossier entier}

\begin{document}

\section*{Résumé}

\begin{resume}
Quatre pages, datées 1964 de sa main --- ce qui est rare dans ce fonds, où
presque toutes les dates sont des déductions d'archiviste --- et qui portent une
liste numérotée de vingt-huit sujets de thèse.

Ce n'est pas un programme de recherche au sens où l'on entend d'ordinaire ce
mot. C'est une liste de travaux à distribuer : vingt-huit questions qu'un
directeur juge assez mûres pour qu'un jeune mathématicien puisse s'y installer
trois ans, assez ouvertes pour qu'il y trouve quelque chose à lui. Grothendieck
prend soin de dire lesquelles sont dures --- il entoure leur numéro et
l'annonce en tête de page --- et de dire de quoi chacune relève : la cohomologie
pour les dix premières, la théorie des intersections pour cinq d'entre elles,
la théorie de Néron pour trois, les topos pour trois autres.

L'intérêt de ces pages n'est donc pas dans une démonstration : il n'y en a
aucune. Il est dans ce qu'une liste de sujets révèle et qu'un article ne
révèle jamais --- l'état d'un domaine vu par celui qui l'organise, à un moment
donné, avec les proportions qu'il lui donne. On y voit ce qui est tenu pour
acquis (assez pour qu'on puisse bâtir dessus), ce qui est tenu pour accessible
(assez pour qu'on l'attribue à un débutant), et ce qui est tenu pour hors de
portée. On y voit surtout la répartition : combien de forces vers la
cohomologie étale, combien vers les cycles, combien vers ce qui n'a pas encore
de nom.

Et l'on y voit une chose que la liste ne prévoyait pas. Les marges de gauche
sont annotées plus tard, dans une autre encre, et ces annotations disent ce
qu'est devenu chaque sujet : qui l'a pris, ce qui a été abandonné, ce qui a été
résolu. Deux de ces notes sont d'une brièveté saisissante. En face du sujet 22
--- structure des catégories abéliennes à produit tensoriel rigide, « préliminaire
algébrique à la théorie des motifs » --- il a écrit \emph{travail Saavedra en
train}. En face du sujet 25, qui demande si un 0-cycle algébriquement
équivalent à zéro sur une surface est rationnellement équivalent à zéro, il a
écrit \emph{faux, cf Mumford}. La liste enregistre ainsi, sur la même feuille,
la question et la réponse : à quelques années de distance, et sans commentaire.

Ce qui suit reprend les vingt-huit sujets dans l'ordre de la page, en donnant à
chacun le nom sous lequel on le chercherait aujourd'hui.

\keywords{Ogg-Shafarevich formula, tannakian categories, Néron model, rational equivalence of zero-cycles, standard conjectures, motives}
\end{resume}

\pagerange{1}{1}
\section*{Ce que la page annonce d'elle-même}

Trois indications ouvrent la liste et gouvernent sa lecture.

D'abord la date, 1964, encadrée en haut à gauche. Elle est de sa main, et c'est
le seul renseignement de ce dossier qui ne soit pas une reconstitution.

Ensuite une convention : \emph{les numéros cerclés sont sans doute difficiles}.
Quinze des vingt-huit sujets portent ce cercle.\footnote{Le décompte est le
nôtre : la page ne totalise pas. Sur quelques numéros le cercle est tracé
légèrement et le relevé n'est pas absolument sûr ; la transcription les donne
tels qu'ils se lisent.} La difficulté est donc déclarée, non pas laissée à
deviner --- ce qui, dans une liste destinée à être distribuée, est une
information de premier ordre.

Enfin une classification, donnée entre parenthèses, qui répartit les sujets par
domaine :

\begin{itemize}
\item cohomologie : 1 à 10, 20, 21, 26, 27, 28 ;
\item théorie des intersections : 10, 19, 23, 24, 25 ;
\item théorie de Néron : 11, 12, 13 ;
\item théorie des topos : 14, 15, 22 ;
\item représentabilité : 17, 18.
\end{itemize}

Le sujet 10 figure dans deux classes, ce qui n'est pas une inadvertance : il
porte sur Riemann--Roch, qui est précisément où la cohomologie et les
intersections se rencontrent. Une note marginale, isolée en haut de page,
ajoute une cinquième entrée : \emph{relation avec cohomologie cristalline}.

\pagerange{1}{1}
\section*{Cohomologie (sujets 1 à 10)}

Le premier bloc porte sur ce que la topologie fppf et la topologie étale
peuvent recevoir de la cohomologie classique.\footnote{Le sujet 1 est écrit très
vite et un mot y est illisible ; la transcription le signale. Le contexte ---
« formalisme général », suivi immédiatement de « dualité (locale, globale) » ---
ne laisse guère de doute sur ce dont il s'agit, mais la restitution est la
nôtre.} Les sujets 1 à 3 sont explicitement liés : le 2 est marqué
\emph{subordonné à 1) et 11)}, et le 3 \emph{subordonné à 2), donc à 1) et 11)}.
Le sujet 2 est d'ailleurs le seul de la liste dont la ligne soit restée vide :
le numéro et la dépendance sont écrits, l'énoncé ne l'est pas.

Les sujets 3 à 6 tournent autour d'une même formule, celle qui calcule la
caractéristique d'Euler--Poincaré d'un faisceau constructible sur une courbe en
fonction de son rang et de ses conducteurs de Swan --- ce qu'on appelle
aujourd'hui la formule de Grothendieck--Ogg--Shafarevich.\footnote{La page écrit
« Ogg--Chafarevitch », transcription française du nom russe. Elle ne nomme pas
Grothendieck lui-même dans la formule, ce qui n'est pas un oubli : le nom à
trois termes est postérieur et n'est pas le sien.} Le sujet 3 en demande la
partie $p$-primaire, celle que la formule laisse hors de portée en
caractéristique $p$ ; le 4 en demande l'analogue en cohomologie étale à
coefficients discrets, avec les représentations d'Artin supérieures ; le 5 une
variante pour les variétés abéliennes munies d'un anneau d'opérateurs.

Le sujet 6 demande des formules de Lefschetz explicites, et le 7 --- porté
difficile --- la cohomologie des schémas de type fini sur $\mathbb{Z}$ avec des
théorèmes de dualité globale « satisfaisants, améliorant Artin ». C'est la
dualité arithmétique globale ; le jugement porté sur l'état de l'art est de la
page.

Les sujets 8 et 9 portent sur le formalisme de dualité du point de vue des
opérateurs différentiels et des coefficients à connexion intégrable, avec un
lien explicite aux « cristaux » --- le mot est de lui, dans son sens de 1964 ---
puis sur la dualité analytique. Le sujet 10 est Riemann--Roch et les problèmes
connexes pour les faisceaux de modules, avec un renvoi qui se lit
\emph{SGA 6}.\footnote{Le sigle est peu lisible et la transcription le marque
comme incertain. S'il est bien SGA 6, il est remarquable : le séminaire
Riemann--Roch a eu lieu en 1966--1967, soit après la date portée sur ces pages,
et la mention serait alors une annotation postérieure comme celles des marges.
On ne tranche pas.}

\pagerange{2}{2}
\section*{Théorie de Néron et théorie des topos (sujets 11 à 16)}

Le sujet 11 demande le foncteur de Néron relatif à un anneau de valuation
discrète, et le 12 « la mise au point, et extensions diverses, de la théorie des
modèles de Néron ». La marge porte : \emph{commencé par Raynaud}.

Le sujet 13 --- difficile --- porte sur des théorèmes asymptotiques pour les
foncteurs de Greenberg au-dessus d'un anneau local noethérien complet ou
hensélien. Sa marge est la plus intéressante de cette page : elle renvoie aux
\emph{résultats d'approximation d'Artin}.\footnote{Le premier mot de cette note
est illisible et la transcription le laisse en blanc ; on ne sait donc pas si la
note dit que le sujet est simplifié par ces résultats, ou dépassé par eux, ou
seulement à reprendre à leur lumière. Le théorème d'approximation d'Artin date
de 1968-1969, ce qui donne à cette annotation un terminus post quem et confirme
que les marges ont été écrites bien après la liste.}

Les sujets 14 et 15 sont les deux sujets « topos » proprement dits. Le 14
demande la structure des topos --- questions d'existence de foncteurs fibres ---
et celle des champs, avec les champs représentables. Le 15, porté difficile,
demande une théorie du type d'homotopie des topos : une définition des
$\pi_{i}$, un théorème de Hurewicz, des suites exactes d'homotopie, un théorème
de comparaison. Sa marge dit : \emph{fait partiellement par M. Artin. Réflexions
de Quillen --- en train}.\footnote{Le nom est abrégé et la transcription le donne
comme incertain ; il peut se lire « M. Artin » ou « Mme Artin ». Le sujet lui-même
--- l'homotopie étale --- est celui sur lequel Michael Artin et Barry Mazur
travaillaient à cette période, ce qui rend la première lecture plus vraisemblable,
mais la page ne le dit pas et cette édition ne le décide pas.}

Le sujet 16 demande une théorie de spécialisation du groupe fondamental pour
une famille de courbes, et sa marge signale un \emph{contre-exemple de Artin}.

\pagerange{2}{2}
\section*{Représentabilité et amplitude (sujets 17 à 19)}

Le sujet 17 demande les cas de représentabilité et de non-représentabilité du
foncteur de Picard $\underline{\mathrm{Pic}}_{X/S}$ pour $X$ de dimension
relative 1 sur $S$. Sa marge porte deux mots : \emph{résolu par Raynaud}.

Le sujet 18, difficile, porte sur les passages au quotient dans des cas
arbitraires et sur l'effectivité des données de descente plates --- deux
questions qui sont, pour la théorie des champs algébriques, la difficulté
centrale. Le 19 demande une théorie de l'amplitude pour les cycles de dimension
quelconque.

\pagerange{3}{3}
\section*{Un encadré : « Pour mémoire »}

Au milieu de la troisième page, un encadré interrompt la liste et note ce qui,
n'étant pas un sujet de thèse, doit tout de même figurer. Deux entrées y sont
marquées d'un signe et reliées par une accolade d'où partent deux implications :

\begin{itemize}
\item la résolution des singularités et la pureté cohomologique
entraîneraient, la première, que la dualité locale soit une bidualité pour la
topologie étale, et la seconde des théorèmes de finitude pour les morphismes de
type fini ;
\item viennent ensuite, sans implication attachée, le changement de base
régulier en cohomologie, la dimension cohomologique des morphismes affines, des
théorèmes de finitude « genre Grauert », et des théories de type
rigide-analytique.
\end{itemize}

Cet encadré est le seul endroit de la liste où Grothendieck écrive des
implications plutôt que des questions : ce sont les énoncés dont il sait ce
qu'ils donneraient, et dont il ne sait pas s'ils sont vrais.

\pagerange{3}{4}
\section*{Théorie des intersections (sujets 23 à 25)}

Un titre --- \emph{Questions dans SGA 1962} --- ouvre le dernier bloc.

Le sujet 23 demande une théorie de Chow, c'est-à-dire une théorie de
l'équivalence rationnelle des cycles, pour les schémas munis d'un faisceau
inversible ample, avec cette exigence précise : que l'image inverse par
$f : X \to Y$ soit définie sous la seule hypothèse que $Y$ soit
régulier.\footnote{Sur la page, le mot « réguliers » qualifiant les schémas est
barré, et c'est la régularité de $Y$ seule qui subsiste dans l'exigence. La
distinction est celle qui gouverne aujourd'hui la construction des
\emph{pullbacks} raffinés en théorie de l'intersection.}

Le sujet 24, difficile, demande si sur une variété projective lisse connexe sur
un corps algébriquement clos un 1-cycle numériquement équivalent à zéro l'est
aussi pour une équivalence plus fine, et en quelle dimension l'équivalence
numérique coïncide avec l'équivalence homologique.\footnote{Les deux dernières
lignes de cette page courent jusqu'au bord du cliché et les dimensions
annoncées n'y sont plus lisibles ; la transcription les laisse en blanc et
l'énoncé donné ici est en conséquence plus général que celui de la page. La
marge, également coupée, porte : \emph{réponse négative Griffiths pour}\ldots}
La question de savoir où l'équivalence numérique et l'équivalence homologique
se confondent est l'une des conjectures standard.

Le sujet 25, difficile, demande si le produit de deux cycles algébriquement
équivalents à zéro est rationnellement équivalent à zéro, puis, en particulier,
si un 0-cycle algébriquement équivalent à zéro sur une surface est
rationnellement équivalent à zéro. La marge porte trois mots : \emph{faux, cf
Mumford}.\footnote{L'annotation ne donne ni titre ni date. Elle est conforme à
ce qui est établi depuis --- la réponse à cette question est négative, et le
groupe de Chow des 0-cycles d'une surface de genre géométrique non nul n'est pas
de dimension finie --- mais l'identification précise de la référence est la
nôtre, non celle de la page.}

\pagerange{4}{4}
\section*{Les quatre derniers (sujets 26 à 28)}

Le sujet 26 demande une théorie de Rosenlicht relative, sur base quelconque ;
le 27, les adèles des schémas de dimension quelconque et les jacobiennes
locales. Le 28, difficile, ferme la liste :

\begin{quote}
Construction d'une théorie abstraite des motifs sur préschémas de type fini sur
$\mathbb{Z}$ (construction catégorique générale).
\end{quote}

C'est, avec le sujet 22, le second endroit où la liste nomme les motifs. Les
deux se répondent : le 22 demande l'algèbre --- la structure des catégories
abéliennes à produit tensoriel rigide, en termes de représentations de groupes
proalgébriques et de gerbes, explicitement présenté comme \emph{préliminaire
algébrique à la théorie des motifs} --- et le 28 demande la construction
elle-même. Le premier porte en marge \emph{travail Saavedra en
train}.\footnote{C'est le programme tannakien. Le rapprochement entre l'énoncé
du sujet 22 et ce qu'on appelle aujourd'hui catégorie tannakienne est immédiat
et n'a rien de conjectural : « catégorie abélienne à produit tensoriel rigide »
et « représentations d'un groupe proalgébrique » sont les deux termes de
l'équivalence. Ce que l'annotation marginale établit, en revanche, est plus
précis et plus rare : que le sujet était, à la date où elle a été portée,
confié et en cours.}

Le reste de la page est vierge.

\end{document}
